\section{Discussion and Limitations}
OrnaScout addresses a systems question: how should a limited multimodal budget be allocated and accepted when source images are much larger than a compact VLM's visual context? The answer need not be one universally dominant selector. Detector ranking, spatial diversity, and negative evidence optimize different objectives; the useful result is the measured operating frontier and the conditions under which each policy improves source-level utility or selective risk.

The VLM stage is similarly treated as an empirical hypothesis. Presence, count, density, and coverage can often be derived from detections, so detector-only and deterministic aggregation are necessary controls. Structured VLM inference is justified only if it adds source-level utility, auditable evidence, or better selective behavior at a transparent computational and review cost. Otherwise, the simpler detector system is the appropriate conclusion.

Several limitations bound the study. The corpus contains two ornamental crop domains and is imbalanced in source count, annotation density, and object scale. One-class boxes support inventory-oriented outcomes but not health, stress, quality, mortality, yield, or treatment recommendations. The selector inherits detector misses; exploration can expose overlooked regions but cannot guarantee recovery. Empirical calibration can fail after severe target shift and does not provide distribution-free risk control. Source-level review also incurs human cost, which must be reported alongside machine latency. Finally, private imagery can support a controlled systems study but is not a dataset contribution unless release rights, documentation, and access are established.

The intended use is decision support. Accepted records summarize visual evidence under a declared protocol; review referrals identify sources for human inspection. Neither path replaces agronomic judgment.

\section{Conclusion}
This revision formulates high-resolution UAV nursery analysis as fixed-budget, source-level evidence allocation with selective review. The completed paper will establish whether structured VLM evidence adds value beyond detector-only aggregation, quantify directed crop transfer, and report the utility--risk--compute frontier of equal-budget policies. \pending{Replace with the final artifact-verified finding and bounded implication before submission.}
